% Author continuation. Independent reconstruction is recorded separately.
% Native dependencies: 315, 327--329 and 335.
\section{A two-block selected multiplier and linear-width roof payment}
\label{sec:sc-selected-multiplier}

The complete selected endpoint kernel can be inserted in both actual
blocks used by the incidence--inverse-gradient argument.  This improves
the physical interior estimate without replacing an outgoing
roof-boundary estimate by an endpoint-momentum estimate.  We retain the
original cap, physical clock, section, exact labels and conditioning.

Write $C_{\rm ce}$ for the particular constant $C_{\rm end}$ in
\eqref{eq:ce-constants}, to distinguish it from other endpoint constants.
For the actual event $E=E_{j,\delta,R}$ in
\eqref{eq:lc-selected-roof}, set
\begin{align}
 K_{\rm sc}&=3C_{\rm ce}/\pi^{2/3},&
 \theta&=\theta_{\rm r},&q_0&=47/53,\notag\\
 d&=\min(j,m-1-j),& L&=\max(j,m-1-j),\notag\\
 \alpha&=\delta+\theta^L,&\beta&=\delta+\theta^d,
                                                   \label{eq:sc-factors}\\
 L_i&=\begin{cases}
       \max(j,i-1-j),&j<i,\\
       \max(j-i,m-1-j),&i\le j,
      \end{cases}
 &\beta_i&=\delta+\theta^{L_i}\quad(1\le i<m).
                                             \label{eq:sc-block-factors}
\end{align}
In both cases $L_i\ge d$, hence $\beta_i\le\beta$.
Keep $r=(p_0^2+p_m^2)^{1/2}$ and the unchanged multiplier
$Z_m=I_m/(Rr)$, with
$I_m=1+\sum_{i=1}^{m-1}q_0^{d_i}/c_i$.

\begin{lemma}[Two selected physical block kernels]
\label{lem:sc-intersection}
For every $R\in[R_-,R_+]$, $s,u>0$, integer $m\ge1$,
$0\le j<m$ and $0<\delta\le\delta_*$, the unconditioned collision measure satisfies
\begin{align}
 \nu(E\cap\{r<s\})
   &\le (C_{\rm ce}/4)\alpha s^2,                 \label{eq:sc-disk}\\
 \nu(E\cap\{r<s,c_i<u\})
   &\le (C_{\rm ce}/\pi)\beta_i s u^2
    \le (C_{\rm ce}/\pi)\beta s u^2,\quad1\le i<m.
                                           \label{eq:sc-intersection}
\end{align}
Every physical word and witness is already summed.  No section factor
or independence assumption enters these inequalities.
\end{lemma}
\begin{proof}
Integrate \eqref{eq:ce-kernel} in the actual initial momentum with
$d\nu=d\alpha_0\,dp_0/(4\pi)$.  The selected $(p_0,p_m)$ density is
bounded by $C_{\rm ce}\alpha/(4\pi)$.  A disk of radius $s$ has area
at most $\pi s^2$, proving \eqref{eq:sc-disk}.

If $j<i$, use the actual block from collision $0$ to collision $i$.
Its length is $i$ and its selected flight is $j$, giving the first
value of $L_i$.  Drop later-record restrictions only on this positive
upper side.  Since $r<s$ implies $|p_0|<s$, the relevant rectangle
in $(p_0,p_i)$ has area at most $4su^2$: the first interval has
length at most $2s$, and the set $c_i<u$ has momentum length
$2(1-\sqrt{1-u^2})\le2u^2$ for $u\le1$ and at most $2\le2u^2$
for $u>1$.

If $i\le j$, invariance starts the actual block at collision $i$.
Its length is $m-i\ge1$ and its selected flight is $j-i$, giving
the second value of $L_i$.  Retain $|p_m|<s$ and the same grazing
set for $p_i$, now in $(p_i,p_m)$.  Apply the same rectangle bound
to the selected density $C_{\rm ce}\beta_i/(4\pi)$.
In particular $i=j$ uses the first-flight case of
Theorem~\ref{thm:ce-kernel}.  Both block constructions retain the same
outgoing flight, actual first hit and unvisited witness.  The
rectangle area under a proved joint-density bound is not a product
of physical probabilities.  This proves \eqref{eq:sc-intersection}.
\end{proof}

\begin{theorem}[The complete selected singular multiplier]
\label{thm:sc-multiplier}
For every Borel $0\le h\le\one_E$, put $a_h=\int h\,d\nu$.
For $1\le i<m$ and $t>0$,
\begin{align}
 \nu\{x\in E:(rc_i)^{-1}>t\}
 &\le K_{\rm sc}\alpha^{1/3}\beta_i^{2/3}t^{-4/3},
                                             \label{eq:sc-mixed-tail}\\
 \int h/(rc_i)\,d\nu
 &\le4K_{\rm sc}^{3/4}\alpha^{1/4}\beta_i^{1/2}a_h^{1/4}.
                                       \label{eq:sc-mixed-functional}
\end{align}
Consequently,
\begin{align}
 \int hZ_m\,d\nu
 &\le R_-^{-1}\left\{\sqrt{C_{\rm ce}}\alpha^{1/2}a_h^{1/2}
 +4K_{\rm sc}^{3/4}\alpha^{1/4}a_h^{1/4}
                  \sum_{i=1}^{m-1}q_0^{d_i}\beta_i^{1/2}\right\}
                                            \label{eq:sc-full-refined}\\
 &\le R_-^{-1}\left\{\sqrt{C_{\rm ce}}\alpha^{1/2}a_h^{1/2}
 +\frac{188}{3}K_{\rm sc}^{3/4}
                       \alpha^{1/4}\beta^{1/2}a_h^{1/4}\right\}.
                                                     \label{eq:sc-full}
\end{align}
The constants are uniform in all original counts, radii, indices,
labels and roofs.  At $m=1$ the incidence sum is empty.
\end{theorem}
\begin{proof}
Fix $i$ and write $A=C_{\rm ce}\alpha/4$,
$B=C_{\rm ce}\beta_i/\pi$.  For any $b>0$, the part $r<b$ costs
$Ab^2$.  On the annulus $2^kb\le r<2^{k+1}b$, the selected event
$(rc_i)^{-1}>t$ costs at most $2B(t^2b)^{-1}2^{-k}$ by
\eqref{eq:sc-intersection}.  Summation gives
\[
 \nu\{x\in E:(rc_i)^{-1}>t\}\le Ab^2+4B/(t^2b).
\]
The minimizer is $b=(2B/A)^{1/3}t^{-2/3}$, with value
\[
 3\,2^{2/3}A^{1/3}B^{2/3}t^{-4/3}
      =K_{\rm sc}\alpha^{1/3}\beta_i^{2/3}t^{-4/3}.
\]
This value is the exact minimum of the displayed annular upper bound
over $b>0$, not an assertion of the best constant in the theorem.
This is an analytical cutoff in a finite-count integral, not a
changed physical event or roof window.  Layer cake integrates
$\min\{a_h,K_{\rm sc}\alpha^{1/3}\beta_i^{2/3}t^{-4/3}\}$
to give \eqref{eq:sc-mixed-functional}.  At $a_h=0$ use bounded
truncation.  The disk inequality similarly gives
\[
 \int h/r\,d\nu
 \le2\sqrt{C_{\rm ce}\alpha/4}\,a_h^{1/2}
 =\sqrt{C_{\rm ce}}\alpha^{1/2}a_h^{1/2}.
\]
Expand the actual $I_m$ to get \eqref{eq:sc-full-refined}; then use
$\beta_i\le\beta$ and the inherited deterministic sum
$\sum_{i=1}^{m-1}q_0^{d_i}\le47/3$.  All integrals concern one orbit
after the complete kernel sums, not products of independent blocks.
\end{proof}

\begin{remark}[Mixed Lorentz endpoint and actual scope]
\label{rem:sc-endpoint}
For $0<p<4/3$, the identical layer-cake calculation gives
\[
 \int h(rc_i)^{-p}\,d\nu
 \le\frac4{4-3p}
  (K_{\rm sc}\alpha^{1/3}\beta_i^{2/3})^{3p/4}
                                      a_h^{1-3p/4}.
\]
At $p=4/3$ only the weak tail \eqref{eq:sc-mixed-tail} is proved,
not a finite strong moment.  The exact $p=1$ calculation already
gives the linear-width result below; no new physical norm is required.
\end{remark}

\subsection{The roof height with outgoing boundaries retained}
Define
\begin{equation}\label{eq:sc-interior-budget}
 \mathcal I(\delta,\alpha,\beta)
 =\frac{\sqrt{C_{\rm ce}}}{R_-^{3/2}}(\delta\alpha)^{1/2}
  +\frac{188K_{\rm sc}^{3/4}}{3R_-^{5/4}}
                \delta^{1/4}\alpha^{1/4}\beta^{1/2}.
\end{equation}
Keep exactly $C_{\rm int}$ from \eqref{eq:rbv-constants} and
$C_{\rm side},C_\delta,C_{\rm cap,loc}$ from
\eqref{eq:lc-boundary-constants}.

\begin{theorem}[Selected original-roof height with two-block gain]
\label{thm:sc-roof}
The same selected density $p_{m,R;j,\delta}$ as in
Theorem~\ref{thm:lc-selected-roof-height} satisfies
\begin{equation}\label{eq:sc-roof}
\begin{split}
 \|p_{m,R;j,\delta}\|_\infty
 \le{}& (C_{\rm int}/c_*)\mathcal I(\delta,\alpha,\beta)
           +(C_{\rm side}+C_{\rm cap,loc})\delta\\
      &+(C_{\rm side}+C_\delta)(\theta^j+\theta^{m-1-j}).
\end{split}
\end{equation}
If $\theta^d\le\delta\le\delta_*$, then
\begin{equation}\label{eq:sc-linear-roof}
                    \|p_{m,R;j,\delta}\|_\infty\le C_{\rm lin}\delta,
\end{equation}
with fixed admissible constants
\begin{align}
 J_{\rm lin}
   &=\frac{\sqrt{2C_{\rm ce}}}{R_-^{3/2}}
     +\frac{188\,2^{3/4}K_{\rm sc}^{3/4}}{3R_-^{5/4}},
                                        \label{eq:sc-linear-constants}\\
 C_{\rm lin}
   &=(C_{\rm int}/c_*)J_{\rm lin}
       +3C_{\rm side}+2C_\delta+C_{\rm cap,loc}.\notag
\end{align}
All original exact-event and Borel selectors $0\le H\le1$ remain
available by positive restriction.
\end{theorem}
\begin{proof}
Invariance and \eqref{eq:cf-exact-mean} give
$\nu(E)\le\delta/R\le\delta/R_-$.  Thus \eqref{eq:sc-full} with
$h=\one_E$ bounds $\int_E Z_m\,d\nu$ by
$\mathcal I(\delta,\alpha,\beta)$.  The selected interior calculation
from \eqref{eq:rbv-interior-pointwise} is consequently bounded by
$C_{\rm int}\mathcal I/c_*$.  Only this interior estimate is replaced.

The original selected old boundaries, new width-$\delta$ boundaries
and residual endpoint caps still have exactly the budgets
\eqref{eq:lc-old-boundary}, \eqref{eq:lc-new-boundary} and
Lemma~\ref{lem:lc-selected-cap}.  Apply the same physical
right-endpoint inequality \eqref{eq:lc-right-endpoint} to the complete
selected intervals.  Visited grazing retains its actual zero trace;
every positive-visited critical endpoint remains incoming, and no
critical coefficient is discarded.  Sum these unchanged outgoing
budgets and the new interior estimate to prove \eqref{eq:sc-roof}.
Absolute continuity and positive-selector statements follow from the
same cap domination.  The selector is not differentiated, and the
unique conditioning factor is $c_*^{-1}$ in the interior and is
already contained in the boundary constants.

If $\theta^d\le\delta$, then $\alpha\le\beta\le2\delta$ and
$\theta^j+\theta^{m-1-j}\le2\delta$.  The interior budget is at most
$J_{\rm lin}\delta$, and the boundary terms are at most
$(3C_{\rm side}+2C_\delta+C_{\rm cap,loc})\delta$.  This proves
\eqref{eq:sc-linear-roof}.
\end{proof}

\begin{corollary}[Linear selected moments and relative tail control]
\label{cor:sc-linear-moments}
Assume $\theta^d\le\delta\le\delta_*$ and put
$W_{\rm lin}=R_-^{1/4}J_{\rm lin}$.  Then for every Borel
$0\le h\le\one_E$, with $a_h=\int h\,d\nu$,
\begin{align}
 \int hZ_m\,d\nu
      &\le W_{\rm lin}\delta^{3/4}a_h^{1/4},
                                             \label{eq:sc-relative-ui}\\
 \nu(E\cap\{Z_m>t\})
      &\le W_{\rm lin}^{4/3}\delta\,t^{-4/3}
                         \qquad(t>0),          \label{eq:sc-full-tail}\\
 \int_E Z_m^p\,d\nu
      &\le\frac4{4-3p}R_-^{p-1}J_{\rm lin}^p\,\delta
                         \qquad(1<p<4/3),     \label{eq:sc-linear-moments}\\
 \int_{E\cap\{Z_m>T\}}Z_m\,d\nu
      &\le4W_{\rm lin}^{4/3}\delta\,T^{-1/3}
                         \qquad(T>0).          \label{eq:sc-relative-tail}
\end{align}
For the unchanged conditional measure $\nu_R^*$, and with any
original additional positive selector, the moment and tail integral
upper bounds are multiplied by $c_*^{-1}$ exactly once.  These are
estimates on the original selected source, not a replacement of its
conditioning event by a clearance-conditioned probability law.
\end{corollary}
\begin{proof}
Use $\alpha,\beta\le2\delta$ in \eqref{eq:sc-full}.
Since $a_h\le\nu(E)\le\delta/R_-$, its first term is bounded by
\[
 \frac{\sqrt{2C_{\rm ce}}}{R_-}
       \delta^{1/2}a_h^{1/2}
 \le\frac{\sqrt{2C_{\rm ce}}}{R_-^{5/4}}
       \delta^{3/4}a_h^{1/4}.
\]
The other term already has this form, with coefficient
$(188/3)2^{3/4}K_{\rm sc}^{3/4}/R_-$.  The sum of these
coefficients is exactly $W_{\rm lin}$, proving
\eqref{eq:sc-relative-ui}.  Set
$h=\one_{E\cap\{Z_m>t\}}$; if its mass $a_h$ is positive,
$ta_h\le W_{\rm lin}\delta^{3/4}a_h^{1/4}$ gives
\eqref{eq:sc-full-tail}.  The zero-mass case is immediate.

Combine that tail with $\nu(E)\le\delta/R_-$ and integrate
$p t^{p-1}\min\{\delta/R_-,W_{\rm lin}^{4/3}\delta t^{-4/3}\}$.
Its integral is
\[
 \frac4{4-3p}
       (\delta/R_-)^{1-3p/4}
       (W_{\rm lin}^{4/3}\delta)^{3p/4}
 =\frac4{4-3p}R_-^{p-1}J_{\rm lin}^p\delta.
\]
For \eqref{eq:sc-relative-tail}, use
$\int_{E\cap\{Z_m>T\}}Z_m
 =T\nu(E\cap\{Z_m>T\})+\int_T^\infty\nu(E\cap\{Z_m>t\})\,dt$.
Finally $d\nu_R^*=\one_{Y_R^*}d\nu/c_*$; positive restrictions
and this identity prove the stated conditional versions without any
division by the selected mass.
\end{proof}

\subsection{An enlarged paid flight range at the original normalization}
Retain the original widths
$\delta_{j,B}=2\varepsilon(B)q_0^{d'_j/4}$ and the actual union
$\mathcal E^{\rm clr}_{B,m}(D)$ in
\eqref{eq:lc-original-depths}--\eqref{eq:lc-deep-event}.  Put
\begin{equation}\label{eq:sc-ratio}
                         \rho=\theta/q_0^{1/4}<1.
\end{equation}
Indeed \eqref{eq:lc-theta} gives $\theta<53/59$, and
$(53/59)^4<47/53=q_0$.

\begin{theorem}[Linear-width deep payment in the original order]
\label{thm:sc-deep-payment}
At every finite integer $m\ge1$, original outer $B$, and integer
$D\ge0$ satisfying
\begin{equation}\label{eq:sc-depth-conditions}
 2\varepsilon(B)q_0^{D/4}\le\delta_*,
                 \qquad \rho^D\le2\varepsilon(B),
\end{equation}
the same original exact-event pushforward and persistent subsources
of Theorem~\ref{thm:lc-deep-payment} have density bounded by
\begin{equation}\label{eq:sc-deep-finite}
 Q^{\rm lin}_{\varepsilon(B)}(D)
       =\frac{4C_{\rm lin}\varepsilon(B)}{1-q_0^{1/4}}q_0^{D/4}.
\end{equation}
For fixed $\eta>0$, put
\begin{equation}\label{eq:sc-log-depth}
 D_m^{\rm lin}
   =\left\lceil\frac{4(2+\eta)\log(m+1)}{|\log q_0|}\right\rceil,
 \qquad
 A_B^{\rm lin}=\frac{4C_{\rm lin}\varepsilon(B)}{1-q_0^{1/4}}.
\end{equation}
At each fixed original outer $B$ with $\varepsilon(B)>0$, all
sufficiently large counts satisfy
\begin{equation}\label{eq:sc-normalized-rate}
 \sup_{R,n,k}m^2
 \|\mathsf U^{\rm per}_{B,m}
       [\mathcal E^{\rm clr}_{B,m}(D_m^{\rm lin})]\|_\infty
              \le A_B^{\rm lin}(m+1)^{-\eta}.
\end{equation}
No physical roof or comparison window, label or original guard
width is changed.
\end{theorem}
\begin{proof}
At every selected depth $d'_j\ge D$, its width is at most
$\delta_*$, and
$\theta^{d'_j}/\delta_{j,B}=\rho^{d'_j}/(2\varepsilon(B))\le1$.
The linear-height theorem therefore applies at each selected flight.
There are at most two indices of each depth, and hence
\[
 \sum_{d'_j\ge D}C_{\rm lin}\delta_{j,B}
 \le \frac{4C_{\rm lin}\varepsilon(B)}{1-q_0^{1/4}}q_0^{D/4}.
\]
Positivity proves \eqref{eq:sc-deep-finite} after every actual word,
mark and impact sign has been included, before the radius and label
suprema.  The elementary persistent domination
\eqref{eq:v111-domination}, with the unchanged target-roof selector,
proves the same estimate for the persistent source and its positive
subsources, exactly as in Theorem~\ref{thm:lc-deep-payment}.

For each fixed outer $B$, both conditions in
\eqref{eq:sc-depth-conditions} eventually hold since
$D_m^{\rm lin}\to\infty$, $q_0,\rho<1$ and
$0<\varepsilon(B)<\infty$.  Also
$q_0^{D_m^{\rm lin}/4}\le(m+1)^{-(2+\eta)}$.
Multiply by $m^2\le(m+1)^2$ to obtain
\eqref{eq:sc-normalized-rate}.  This is the original collision limit
at fixed outer $B$, not a collision-dependent choice of an outer
parameter.  Since $D_m^{\rm lin}=O(\log m)=o(m)$, the paid range is
nonempty and contains $m-O(\log m)$ actual flights for large counts.
\end{proof}

\paragraph{Gain and limitation.}
The new argument uses the selected kernel on two specified blocks of
one physical record.  It proves linear dependence on the actual width
when $\theta^{d'_j}\le\delta$, with all outgoing charges retained.
For the unchanged guard, the leading logarithmic depth coefficient
decreases from $16(2+\eta)/|\log q_0|$ in
\eqref{eq:lc-log-depth} to $4(2+\eta)/|\log q_0|$.
Here $\theta<q_0^{1/4}<q_0^{1/16}$, so the old $q_{\rm c}$ is
indeed exactly $q_0^{1/16}$.

This is a stronger bound on an enlarged class of actual original
sources, not a renaming of an unpaid source.  At a shallow fixed
flight, $\theta^{d'_j}\le\delta_{j,B}$ need not hold.  The outgoing
terms remain $\theta^j+\theta^{m-1-j}$, not
$\theta^{\max(j,m-1-j)}$.  Different physical sheets at a common roof
have different endpoint momenta, and outgoing boundaries have their
own turning measures; a common-endpoint kernel cannot simply replace
their bound.  The present argument does not pay the remaining shallow
source, establish comparison-pool abundance, or prove the full
normalized local derivative charge.  The original raw local-limit
closure and its separate spectral, arithmetic, conditioning and inverse
dependencies remain unresolved.
